\documentclass[11pt,a4paper]{article}

\usepackage[utf8]{inputenc}
\usepackage[T1]{fontenc}
\usepackage[serbian]{babel}
\usepackage{lmodern}
\usepackage{geometry}
\usepackage{setspace}
\usepackage{amsmath,amssymb}
\usepackage{graphicx}
\usepackage{booktabs}
\usepackage{hyperref}
\usepackage{enumitem}
\usepackage{csquotes}
\usepackage[
  backend=biber,
  style=authoryear,
  sorting=ynt
]{biblatex}
\addbibresource{references.bib}

\geometry{
  left=25mm,
  right=25mm,
  top=25mm,
  bottom=30mm
}

\onehalfspacing
\setlength{\parindent}{0pt}
\setlength{\parskip}{6pt}

\title{\textbf{Bajesovska optimizacija}}
\author{
Ognjen Nešković 9/2022 \\
Pavle Sekešan 87/2022 \\
\small Projekat u okviru kursa Racunarska inteligencija
}
\date{Februar 2026}

\begin{document}

\maketitle
\thispagestyle{empty}
\newpage

\tableofcontents
\newpage

\section{Uvod}
Bajesovska optimizacija (\textit{Bayesian optimization}) je porodica algoritama za \textit{Black box}\footnote{U ovom kontekstu Black box optimizacija odnosi se na optimizaciju funkcije bez poznavanja ikakvih njenih svojstava osim vrednosti u odabranim tačkama.} optimizaciju funkcija.
Bajesovska optimizacija se primenjuje kao vrlo generalna metoda u slučajevima kada ne možemo pretpostaviti ništa dodatno u vezi funkcije koja se optimizuje - u praksi ova metoda optimizacije se primenjuje kada je funkcija neprekidna, nije diferencijabilna (ili je presporo evaluirati njen izvod), skupa da se evaluira i domen funkcije je nisko-dimenzion $d \leq 20$ \parencite{frazier2018tutorialbayesianoptimization}.
Česta upotreba Bajesovske optimizacije je kao tehnika za nalaženje optimalnih hiperparametara modela koji se koriste u mašinskom učenju --- broj i veličina slojeva neuronske mreže, aktivacione funkcije, maksimalna dubina stabala odlučivanja...

Osnovna ideja algoritama Bajesovske optimizacije je sledeća: Kako funkcija $f$ koja se optimizuje nije poznata tretiraćemo je kao nasumičnu. Postavićemo apriornu raspodelu nad $f$, kada izaberemo novu tačku i evaluiramo funkciju u toj tački ažuriraćemo raspodelu kako bi dobili aposteriornu raspodelu.
Na osnovu nekog broja evaluacija funkcije $f$ i tačaka u kojima je evaluirana se dobija raspodela nad $f$. Ova raspodela nam daje u suštini \textit{surogat} funkciju. Ova \textit{surogat} funkcija je aproksimacija koja je značajno lakša da se evaluira nego $f$. Korisno je napomenuti da ovo nije suštinski za tehnike Bajesovske optimizacije.
Postoje i druge optimizacione metode koje se zasnivaju na pronalaženju i optimizaciji \textit{surogat} funkcija \parencite{booker1999rigorous}.

Algoritmi zasnovani na idejama Bajesovske optimizacije se suštinski razlikuju u tome kako postavljaju raspodelu nad $f$ i kako biraju narednu tačku za evaluaciju na osnovu postavljene raspodele. Metoda koja je implementirana u ovom radu se zasniva na tehnici GP (\textit{Gaussian Process}) regression za određivanje raspodele nad $f$.
Na osnovu raspodele koja je određena nad $f$ potrebno je odabrati sledeću tačku u kojoj će se evaluirati $f$. Suštinski, prilikom odabira naredne tačke potrebno je balansirati koliko se istražuje deo prostora gde je nesigurnost funkcije velika i koliko se eksploatišu delovi funkcije za koje je već poznato da su potencijalno blizu optimalne vrednosti.
U ovom radu je u te svrhe implementiran tzv. princip očekivanog poboljšanja (EI - Expected Improvement).

\subsection{Ograničenja}
Metoda Bajesovske optimizacije koja je opisana i implementirana u okviru ovog rada rešava problem:
$$\min_{x \in A}f(x)$$
Metoda se može primeniti ukoliko je ispunjeno sledeće:
\begin{itemize}
  \item Za neko $x$ je lako proveriti da li pripada domenu $A$. Konkretno u ovoj implementaciji se pretpostavlja da je domen oblika $\{ x \in \mathbb{R}^d | l_i \leq x \leq h_i  \}$.
  \item Funkcija $f$ je neprekidna
\end{itemize}

\section{Bajesovska optimizacija GP-EI}


\subsection{GP Regresija}

\subsection{Akviziciona funkcija, EI}

\section{Implementacija}

\section{Eksperimenti}

\subsection{Postavka}

\subsection{Rezultati}

\subsection{Diskusija}

\section{Zakljucak}

\printbibliography

\end{document}


